% Seção VI — Proposta de Solução
\section{Proposta de Solução}
\label{sec:proposta}

\subsection{Visão Geral da Solução}
\label{sec:visao-geral}

[Descrever a proposta em alto nível.
O que será construído? Qual o diferencial tecnológico em relação às soluções existentes?
Como a solução atende ao problema identificado na Seção~\ref{sec:problema}?]

\subsection{Arquitetura do Sistema}
\label{sec:arquitetura}

% Substituir o \fbox pelo \includegraphics após criar o diagrama em /images
\begin{figure}[h]
  \centering
  % \includegraphics[width=0.9\columnwidth]{images/diagrama_arquitetura}
  \fbox{\parbox{0.85\columnwidth}{\centering\vspace{2cm}
    [PLACEHOLDER: Diagrama de Arquitetura UML ou Game Design Document]
  \vspace{2cm}}}
  \caption{Arquitetura proposta do sistema.}
  \label{fig:arquitetura}
\end{figure}

\subsection{Fluxo Principal}
\label{sec:fluxo}

[Descrever o fluxo principal de uso do sistema passo a passo.
Pode incluir um diagrama de sequência ou de casos de uso como figura adicional.
Referenciar a Figura~\ref{fig:arquitetura} quando pertinente.]

\subsection{Tecnologias e Justificativa de Escolha}
\label{sec:justificativa-tec}

[Para cada tecnologia escolhida, justificar tecnicamente por que ela foi selecionada
em detrimento de alternativas disponíveis no mercado.
Considerar aspectos como: maturidade, suporte da comunidade, licença, desempenho e
compatibilidade com as demais ferramentas da Tabela~\ref{tab:ferramentas}.]
